\documentclass[conference]{IEEEtran}
\IEEEoverridecommandlockouts

\usepackage{cite}
\usepackage{amsmath,amssymb,amsfonts}
\usepackage{algorithmic}
\usepackage{algorithm}
\usepackage{graphicx}
\usepackage{textcomp}
\usepackage{xcolor}
\usepackage{booktabs}
\usepackage{multirow}
\usepackage{hyperref}
\usepackage{bbm}

\def\BibTeX{{\rm B\kern-.05em{\sc i\kern-.025em b}\kern-.08em
    T\kern-.1667em\lower.7ex\hbox{E}\kern-.125emX}}

\begin{document}

\title{PRISM: Pareto-optimal Risk-constrained drIving Style Manifold\\
{\footnotesize A Distributional Multi-Objective Framework for Safe and
Personalized Autonomous Driving}}

\author{
\IEEEauthorblockN{Anonymous Author(s)}
\IEEEauthorblockA{Anonymous Institution}
}

\maketitle

%----------------------------------------------------------
\begin{abstract}
Existing reinforcement learning (RL) planners for autonomous driving
optimize a single scalar reward that encodes one fixed driving style and
treats safety as a tradeable objective. This fails to serve users with
diverse preferences and provides no formal safety guarantee against
catastrophic tail-risk outcomes. We propose PRISM, a framework that
learns a \emph{Distributional Pareto-Optimal} set of driving style
policies---covering comfort-seeking, progress-oriented, disciplined, and
cautious behaviours---while enforcing a hard Conditional Value-at-Risk
(CVaR) safety constraint that is never traded away regardless of style
preference. Building on Distributional Pareto-Optimal Multi-Objective
Reinforcement Learning (DPMORL)~\cite{DPMORL2023}, we optimize over the
entire return distribution of each style objective rather than
expectations alone, producing policies that are robust to return
uncertainty. Safety is enforced via a principled two-tier cost function
separating outcome events from leading indicators, with indicator weights
derived from empirical causal lead times to catastrophic outcomes. We
further propose an \emph{$\alpha$-curriculum}: progressively tightening
the CVaR confidence level from $\alpha\!=\!0.20$ to $\alpha\!=\!0.95$
over training, co-scheduled with the baseline CVaR curve, to avoid
over-conservative early training. We evaluate on the nuPlan closed-loop
benchmark~\cite{nuPlan2024}, the most demanding publicly available
planning benchmark for autonomous driving.
\end{abstract}

\begin{IEEEkeywords}
autonomous driving, multi-objective reinforcement learning,
safe reinforcement learning, CVaR, distributional RL,
driving style personalisation, nuPlan
\end{IEEEkeywords}

%----------------------------------------------------------
\section{Introduction}\label{sec:intro}

Reinforcement learning has emerged as a promising paradigm for
autonomous vehicle (AV) planning~\cite{CaRL2025,PDM2023}. However,
contemporary RL planners share two fundamental limitations.

\textbf{Single driving style.}
Scalar reward weights encode a fixed behavioural trade-off. Different
users have genuinely different driving preferences: some prefer smooth,
comfort-oriented riding; others prefer assertive, progress-seeking
behaviour~\cite{Yusof2016,Vasile2023}. A single policy cannot serve all
users, and retraining per user is impractical. This inflexibility
reduces trust and acceptance of autonomous
vehicles~\cite{Haselberger2024}.

\textbf{Expectation-only optimisation and unsafe safety treatment.}
Current MORL approaches for driving optimise only the \emph{expected}
return of each objective~\cite{Surmann2025}, treating two policies
identically if they share the same expected returns even when their
variance profiles differ fundamentally. Furthermore, existing approaches
either include safety as a tradeable Pareto objective---allowing safety
to be sacrificed for style---or use a fixed penalty that distorts
behavioural objectives.

PRISM addresses both problems simultaneously through two contributions:

\begin{enumerate}
  \item \textbf{Distributional Pareto-Optimal style policies.}
    We learn policies whose \emph{entire return distributions} are
    Pareto-optimal under stochastic dominance~\cite{DPMORL2023},
    distinguishing consistently good policies from merely
    average-good ones.

  \item \textbf{Hard CVaR safety constraint with $\alpha$-curriculum.}
    Safety is a non-negotiable CVaR bound on the tail of the cumulative
    safety cost. A progressive $\alpha$-curriculum avoids over-conservative
    early training while converging to a tight tail guarantee.
\end{enumerate}

%----------------------------------------------------------
\section{Related Work}\label{sec:related}

\subsection{RL-Based Autonomous Driving Planners}
CaRL~\cite{CaRL2025} achieves state-of-the-art results on nuPlan
through a simple route-completion reward with PPO, demonstrating that
simplicity combined with scale outperforms complex shaped rewards.
PDM~\cite{PDM2023} combines rule-based proposals with neural scoring.
PRISM builds on CaRL's observation space and PPO implementation.

\subsection{MORL for Autonomous Driving}
Surmann et al.~\cite{Surmann2025} propose a single-policy MORL approach
for personalised driving with runtime preference adaptation. Their method
uses linear scalarisation of expected returns with no formal safety
guarantee. PRISM differs in three ways: (i)~full return distribution
optimisation; (ii)~safety as a hard CVaR constraint; (iii)~nuPlan
evaluation rather than CARLA.

Abouelazm et al.~\cite{BeyondScalar2026} introduce lexicographic safety
via Quantile Dominance. Like PRISM, they argue safety should not be
traded. However, they produce a single policy under a fixed hierarchy
rather than a diverse safe Pareto front, and evaluate on CARLA.

\subsection{Distributional and Risk-Sensitive RL}
Distributional RL~\cite{C51_2017,QRDQN2018} models full return
distributions. DPMORL~\cite{DPMORL2023} extends this to MORL via
stochastic dominance and proves (Theorem~3) that CVaR-optimal policies
belong to the DPO set---the theoretical basis for our combined
framework. CVaR-constrained RL has been studied in general
settings~\cite{CVaRCPO2024,CPPO2022,TRC2022} but not applied to MORL
autonomous driving. The reward design review~\cite{RewardReview2024}
identifies context-aware progress and tail-risk safety as open problems
that PRISM directly addresses.

\subsection{Safe RL for Autonomous Driving}
Most safe RL approaches use expected cost
constraints~\cite{CPO2017,Tessler2018}. Hu et al.~\cite{LSTC2024}
demonstrate that purely long-term constraints allow vehicles to enter
hazardous states during training, motivating our $\alpha$-curriculum
which progressively tightens the constraint.

%----------------------------------------------------------
\section{Problem Formulation}\label{sec:formulation}

\subsection{Multi-Objective MDP}

\begin{equation}
\mathcal{M}=(\mathcal{S},\mathcal{A},\mathcal{P},\mathcal{P}_{0},
             \mathbf{R},c,\gamma,T)
\end{equation}
$\mathcal{S}$: BEV semantic segmentation (inherited from
CaRL~\cite{CaRL2025}); $\mathcal{A}$: continuous acceleration and
steering; $\mathbf{R}:\mathcal{S}\!\times\!\mathcal{A}\!\to\!\mathbb{R}^{4}$:
style reward vector; $c$: scalar safety cost; $\gamma$: discount factor;
$T$: episode horizon.

\subsection{Style Reward Vector}\label{sec:rewards}

\begin{equation}
\mathbf{r}_{t}=\bigl(r^{\text{com}}_{t},\;r^{\text{prog}}_{t},\;
               r^{\text{lat}}_{t},\;r^{\text{spc}}_{t}\bigr)\in(0,1]^{4}
\end{equation}

The four objectives are motivated by empirical studies identifying three
primary axes of human driving-style preference: longitudinal
assertiveness, lateral discipline, and spacing
behaviour~\cite{Yusof2016,NatDriving2023}. Longitudinal behaviour is
further split into comfort and progress because they capture orthogonal
dimensions of assertiveness.

\subsubsection{Comfort $r^{\text{com}}_{t}$}

\begin{equation}
r^{\text{com}}_{t}=
  \exp\!\Bigl(-\frac{j_{\text{lon}}^{2}(t)+j_{\text{lat}}^{2}(t)}
                    {\sigma_{j}^{2}}\Bigr)
\label{eq:comfort}
\end{equation}
$j_{\text{lon}},j_{\text{lat}}$: longitudinal and lateral jerk.
$\sigma_{j}^{2}$ is the empirical jerk variance from IDM warm-up
rollouts, computed once and cached (Section~\ref{sec:hyperparams}).
The exponential form bounds the reward in $(0,1]$ and is consistent
with the normalization scheme across all four
objectives~\cite{RewardReview2024}.

\subsubsection{Progress $r^{\text{prog}}_{t}$}

\begin{equation}
r^{\text{prog}}_{t}=r^{\text{spd}}_{t}\cdot r^{\text{acc}}_{t}
  \cdot\bigl(0.5+0.5\cdot r^{\text{lane}}_{t}\bigr)
\label{eq:progress}
\end{equation}

The multiplicative structure enforces an AND condition: speed,
acceleration, and lane choice must simultaneously be assertive.
The lane-term floor of $0.5$ prevents temporary suboptimal lane
positions from zeroing out the reward entirely.

\textbf{Speed component:}
\begin{equation}
r^{\text{spd}}_{t}=\exp\!\Bigl(
  -\frac{\max(0,\,v_{\text{des}}(t)-v_{\text{ego}}(t))}
        {\beta\cdot v_{\text{des}}(t)}\Bigr)
\end{equation}
Asymmetric: only penalises being \emph{below} $v_{\text{des}}$; excess
speed is handled by the safety cost. Normalisation by $v_{\text{des}}$
makes the signal scale-invariant across road
types~\cite{RewardReview2024}.

\textbf{Regime-conditioned desired speed}, grounded in
IDM~\cite{IDM1999,IDM25Years2025}:
\begin{equation}
v_{\text{des}}(t)=\begin{cases}
  v_{\text{limit}}(t)     & \text{free flow}\\
  v_{\text{lead}}(t)      & \text{car following}\\
  v_{\text{lane\_avg}}(t) & \text{congested}
\end{cases}
\label{eq:vdesired}
\end{equation}

\textit{Regime detection rules.}
\begin{itemize}
  \item \textbf{Congested}: $v_{\text{lane\_avg}}(t)<0.5\cdot v_{\text{limit}}(t)$,
    where $v_{\text{lane\_avg}}$ is the mean speed of vehicles in the
    same lane within the observation horizon.
  \item \textbf{Car following}: lead vehicle present in ego lane
    within the observation horizon and not congested.
  \item \textbf{Free flow}: no vehicle in ego lane within horizon.
\end{itemize}
Congestion is checked first; free flow is the default fallback.
This context-aware formulation addresses the open problem identified
by~\cite{RewardReview2024}.

\textbf{Acceleration component:}
\begin{equation}
r^{\text{acc}}_{t}=1-\exp\!\Bigl(-\frac{|a_{\text{ego}}(t)|}{\gamma_{a}}\Bigr)
\end{equation}
Magnitude captures both assertive acceleration and decisive
braking~\cite{Surmann2025}. $\gamma_{a}$ cached from warm-up rollouts.

\textbf{Lane choice component:}
\begin{equation}
r^{\text{lane}}_{t}=\frac{\text{lane\_index}(t)}{N(t)-1},\quad
N(t)>1;\quad r^{\text{lane}}_{t}=0.5,\quad N(t)=1
\end{equation}
Leftmost (fastest) lane scores 1; rightmost scores 0. The $N(t)=1$
edge case prevents division by zero on single-lane roads.
Lane choice is a progress signal, not a lateral-discipline
signal~\cite{RewardReview2024}.

\subsubsection{Lateral Discipline $r^{\text{lat}}_{t}$}

\begin{equation}
r^{\text{lat}}_{t}=r^{\text{dev}}_{t}\cdot r^{\text{hdg}}_{t}
\label{eq:lateral}
\end{equation}

\textbf{Lane deviation:}
\begin{equation}
r^{\text{dev}}_{t}=
  \delta_{d}+(1-\delta_{d})\cdot
  \exp\!\Bigl(-\Bigl(\frac{d_{\text{lat}}(t)}{\sigma_{d}}\Bigr)^{2}\Bigr),
\quad \delta_{d}=0.3
\end{equation}
$\sigma_{d}=0.2$\,m from empirical lane-keeping
data~\cite{LKAEmpirical2025}. The floor $\delta_{d}=0.3$ prevents
lane-change manoeuvres from being penalised as
violations~\cite{Haselberger2024}.

\textbf{Heading error:}
\begin{equation}
r^{\text{hdg}}_{t}=
  \exp\!\Bigl(-\Bigl(\frac{|\Delta\psi(t)|}{\phi}\Bigr)^{2}\Bigr)
\end{equation}
$\Delta\psi$: heading angle error relative to lane centreline.
$\phi$ cached from warm-up rollouts.

\subsubsection{Spacing $r^{\text{spc}}_{t}$}

\begin{equation}
r^{\text{spc}}_{t}=
  \delta_{s}+(1-\delta_{s})\cdot
  \Bigl(1-\exp\!\Bigl(-\Bigl(\frac{\max(0,TTC(t))}{\tau}\Bigr)^{2}\Bigr)\Bigr),
\quad\delta_{s}=0.2
\label{eq:spacing}
\end{equation}
\begin{equation}
TTC(t)=\frac{d_{\text{lead}}(t)}{v_{\text{ego}}(t)-v_{\text{lead}}(t)},
\quad v_{\text{ego}}>v_{\text{lead}};\quad
TTC(t)=\infty\;\text{otherwise}
\end{equation}
The floor $\delta_{s}=0.2$ keeps the reward non-degenerate during
assertive following; dangerous low TTC is separately handled in~$c_{t}$.
$\tau$ cached from warm-up rollouts.

\subsection{Hyperparameter Computation}\label{sec:hyperparams}

All scaling parameters
$\Theta=\{\sigma_{j}^{2},\beta,\gamma_{a},\phi,\tau\}$
are computed \emph{once} from $N_{\text{warm}}$ IDM warm-up rollouts
before any policy training begins, stored to disk, and reused across all
$K$ policies and all experiments. This ensures reproducibility and
removes hyperparameter tuning from the training loop.

Specifically, $\sigma_{j}^{2}$ is the empirical variance of jerk;
$\beta$ and $\gamma_{a}$ are set so that a one-standard-deviation
shortfall/excess produces a reward of $e^{-1}\!\approx\!0.37$;
$\phi$ and $\tau$ are calibrated analogously.
The $z_{t}$ normalisation statistics $(\mu_{i},\sigma_{i})$ for each
return dimension are also estimated from warm-up rollouts and updated
with an exponential moving average during training
($\beta_{\text{ema}}=0.01$).

\subsection{Distributional Pareto Optimality}\label{sec:dpo}

Given policy $\pi$, the multi-objective cumulative return is:
\begin{equation}
\mathbf{Z}^{\pi}=\sum_{t=0}^{T}\gamma^{t}\mathbf{r}_{t}\in\mathbb{R}^{4}
\end{equation}
with joint distribution $\mu(\pi)$.

\begin{definition}[Stochastic Dominance~\cite{DPMORL2023}]
$\mu_{1}\!\succ_{\!SD}\!\mu_{2}$ iff $\mu_{1}\!\neq\!\mu_{2}$ and
$\mathbb{E}_{z\sim\mu_{1}}[f(z)]\geq\mathbb{E}_{z\sim\mu_{2}}[f(z)]$
for every non-decreasing $f:\mathbb{R}^{4}\!\to\!\mathbb{R}$.
\end{definition}

\begin{definition}[DPO Policy~\cite{DPMORL2023}]
$\pi$ is Distributional Pareto-Optimal (DPO) if no $\pi'$ satisfies
$\mu(\pi')\succ_{\!SD}\mu(\pi)$.
\end{definition}

DPO is strictly stronger than standard Pareto optimality: it prevents
stochastic dominance across the full return distribution, not just
across expected values. DPMORL Theorem~3~\cite{DPMORL2023} proves that
CVaR-optimal policies are DPO, providing the theoretical basis for
combining the distributional Pareto front with a CVaR constraint.

\subsection{Safety Cost Function}\label{sec:safety}

\begin{equation}
c_{t}=c_{t}^{\text{out}}+\sum_{j}c_{t}^{\text{lead},j}
\label{eq:cost}
\end{equation}

\subsubsection{Outcome Events}

\begin{equation}
c_{t}^{\text{out}}=\sum_{e\in\mathcal{E}}W_{e}\cdot\mathbb{1}_{e}(t)
\end{equation}

\begin{table}[h]
\centering
\caption{Outcome Event Weights (calibrated to nuPlan severity
protocol~\cite{nuPlan2024})}
\begin{tabular}{lc}
\toprule
Event & $W_{e}$\\
\midrule
VRU collision (pedestrian/cyclist)    & 100\\
Wrong-direction driving               & 100\\
Vehicle collision                     & 80\\
Red light violation                   & 80\\
Stop sign violation                   & 70\\
Drivable area violation               & 65\\
Object collision                      & 40\\
\bottomrule
\end{tabular}
\label{tab:weights}
\end{table}

VRU and wrong-direction driving share the highest weight: VRU collisions
have zero tolerance~\cite{nuPlan2024}; wrong-direction driving creates
head-on scenarios with combined closing speeds comparable to the most
severe VRU outcomes~\cite{AV_Collision2022}. Vehicle collision is set to
80 to maintain a meaningful gap from VRU~\cite{AV_Liability2023}.
Drivable area violation is assigned 65 rather than the maximum because
its severity is inherently continuous---minor kerb excursions differ
qualitatively from complete road departures~\cite{nuPlan2024}.

\subsubsection{Leading Indicator Weights}

Leading indicators provide a dense gradient signal during training.
Without them, the agent receives almost no gradient during the majority
of timesteps when no outcome event fires~\cite{LSTC2024}.

For indicator $j$ with outcome set $\mathcal{O}_{j}$:
\begin{equation}
w_{j}^{\text{lead}}=\sum_{i\in\mathcal{O}_{j}}
  \frac{W_{i}}{|\mathcal{O}_{j}|\cdot T_{j\to i}}
\label{eq:indweight}
\end{equation}
$T_{j\to i}$ is the empirically measured mean lead time in timesteps
from indicator $j$ to outcome $i$, estimated from IDM warm-up rollouts.
This ensures a violation firing for its natural lead time produces a
cumulative cost equal to the average weight of outcomes it can cause.

\textit{Severity signals} $s_{j}(t)\in[0,1]$:
\begin{align}
s_{\text{TTC}}(t)&=\max\!\Bigl(0,\frac{\tau_{\text{TTC}}-TTC(t)}
  {\tau_{\text{TTC}}}\Bigr)\cdot\mathbb{1}[TTC(t)<1.5\text{s}]\\
s_{\text{THW}}(t)&=\max\!\Bigl(0,\frac{\tau_{\text{THW}}-THW(t)}
  {\tau_{\text{THW}}}\Bigr)\cdot\mathbb{1}[THW(t)<2\text{s}]\\
s_{\text{spd}}(t)&=\max\!\Bigl(0,\frac{v_{\text{ego}}(t)-v_{\theta}(t)}
  {v_{\theta}(t)}\Bigr)\cdot\mathbb{1}[v_{\text{ego}}>v_{\theta}]\\
s_{\text{blind}}(t)&=\mathbb{1}[|\dot{y}_{\text{ego}}|>\theta_{\text{lc}}]
  \cdot\mathbb{1}[\exists\text{ agent in blind spot}]\\
s_{\text{rl}}(t)&=\mathbb{1}[t_{\text{arrival}}(t)>t_{\text{red}}(t)]
\end{align}
where $v_{\theta}(t)=v_{\text{limit}}(t)(1+\delta_{\text{road}})$ is the
road-type-conditioned speed threshold and
$THW(t)=d_{\text{lead}}(t)/v_{\text{ego}}(t)$.

Per-timestep cost: $c_{t}^{\text{lead},j}=w_{j}^{\text{lead}}\cdot s_{j}(t)$.

\textit{Episode-level cap} for persistent indicators
(TTC, THW, speed):
\begin{equation}
\sum_{t=0}^{T}c_{t}^{\text{lead},j}\leq
  \bar{W}_{j}^{\text{out}}=
  \frac{\sum_{i\in\mathcal{O}_{j}}W_{i}}{|\mathcal{O}_{j}|}
\label{eq:cap}
\end{equation}
Transient indicators (blind spot, predicted red light) whose natural
duration matches their lead time require no cap.

\subsection{CVaR Safety Constraint with $\alpha$-Curriculum}
\label{sec:cvar}

\subsubsection{CVaR Constraint}

\begin{equation}
\text{CVaR}_{\alpha}(C^{\pi})=
  \mathbb{E}\bigl[C^{\pi}\mid C^{\pi}\geq\text{VaR}_{\alpha}(C^{\pi})\bigr]
  \leq\epsilon
\label{eq:cvar}
\end{equation}
$C^{\pi}=\sum_{t}\gamma^{t}c_{t}$ is the trajectory-level cumulative
safety cost, a random variable over rollouts. At $\alpha\!=\!0.95$:
even in the worst $5\%$ of trajectories, the expected safety cost is
bounded by $\epsilon$.

\subsubsection{$\alpha$-Curriculum}

Standard CVaR-constrained RL uses a fixed $\alpha$ throughout
training~\cite{CVaRCPO2024,TRC2022,CPPO2022}. Early in training, a
tight tail constraint ($\alpha\!=\!0.95$) is difficult to satisfy for
an inexperienced policy, causing the Lagrange multiplier $\lambda_{k}$
to spike and crush style optimisation before the policy has learned
basic driving competence.

We propose an $\alpha$-curriculum: linearly increasing $\alpha$ from
$\alpha_{\text{start}}$ to $\alpha_{\text{end}}$ over
$N_{\text{cur}}$ training iterations:

\begin{equation}
\alpha(n)=\alpha_{\text{start}}+(\alpha_{\text{end}}-\alpha_{\text{start}})
  \cdot\min\!\Bigl(1,\frac{n}{N_{\text{cur}}}\Bigr)
\label{eq:curriculum}
\end{equation}
with $\alpha_{\text{start}}\!=\!0.20$, $\alpha_{\text{end}}\!=\!0.95$.

At $\alpha\!=\!0.20$, the constraint covers the worst $80\%$ of
trajectories---providing a broad safety signal early in training
analogous to mean-cost constraints. As $\alpha$ increases, the constraint
progressively focuses on catastrophic tail events.

\textbf{Co-scheduled safety budget.}
To prevent the constraint from hardening abruptly as $\alpha$ increases,
$\epsilon$ is co-scheduled with the CaRL baseline CVaR curve, pre-computed
from $N_{\text{warm}}$ IDM rollouts:
\begin{equation}
\epsilon(n)=\widehat{\text{CVaR}}_{\alpha(n)}^{\text{IDM}}
\label{eq:epsilon}
\end{equation}
This means at every stage the constraint requires policies to be at
least as safe as IDM at the current tail level.

\subsection{Full Optimisation Problem}\label{sec:opt}

Find $K\!=\!5$ DPO policies $\{\pi_{1},\ldots,\pi_{K}\}$ where each
$\pi_{k}$ solves:
\begin{equation}
\max_{\pi_{k}}\;\mathbb{E}_{z\sim\mu(\pi_{k})}\bigl[f_{\theta_{k}}(z)\bigr]
\quad\text{s.t.}\quad
\text{CVaR}_{\alpha(n)}\bigl(C^{\pi_{k}}\bigr)\leq\epsilon(n)
\label{eq:fullopt}
\end{equation}
$f_{\theta_{k}}:\mathbb{R}^{4}\!\to\!\mathbb{R}$ is a non-linear,
non-decreasing utility function from DPMORL Stage~1.

%----------------------------------------------------------
\section{PRISM Training Algorithm}\label{sec:algorithm}

\subsection{Stage 1: Utility Function Generation}

Generate $K$ diverse non-decreasing neural networks
$f_{\theta_{1}},\ldots,f_{\theta_{K}}$ over $[0,1]^{4}$ using
DPMORL's diversity objective~\cite{DPMORL2023}:
\begin{equation}
J(\theta_{i})=
  \alpha_{\text{div}}\cdot J^{\text{val}}(\theta_{i})+
  (1-\alpha_{\text{div}})\cdot J^{\text{grad}}(\theta_{i})
\end{equation}
\begin{align}
J^{\text{val}}(\theta_{i})&=\min_{j\neq i}
  \mathbb{E}_{z\sim U([0,1]^{4})}
  \bigl[(f_{\theta_{i}}(z)-f_{\theta_{j}}(z))^{2}\bigr]\\
J^{\text{grad}}(\theta_{i})&=\min_{j\neq i}
  \mathbb{E}_{z_{1},z_{2}}
  \Bigl[\Bigl(\frac{\Delta f_{\theta_{i}}}{\|z_{2}-z_{1}\|}
             -\frac{\Delta f_{\theta_{j}}}{\|z_{2}-z_{1}\|}\Bigr)^{2}\Bigr]
\end{align}
This produces $K$ utility functions differing in both value and
gradient, ensuring genuine coverage of diverse style preferences.

\subsection{Stage 2: Policy Optimisation}

\subsubsection{State Augmentation}
\begin{equation}
\tilde{s}_{t}=(s_{t},z_{t}),\quad
z_{t}=\sum_{\tau=0}^{t}\gamma^{\tau}\mathbf{r}_{\tau},\quad
z_{t+1}=z_{t}+\gamma^{t}\mathbf{r}_{t}
\end{equation}
$z_{t}$ is normalised using statistics from warm-up rollouts,
updated via EMA during training.

\subsubsection{Scalar Reward via Utility Difference}
\begin{equation}
R_{t}=\gamma^{-t}\bigl[f_{\theta_{k}}(z_{t+1})-f_{\theta_{k}}(z_{t})\bigr]
\label{eq:scalar}
\end{equation}
DPMORL Theorem~2 guarantees the optimal policy of this augmented MDP
equals the optimal policy under $f_{\theta_{k}}$~\cite{DPMORL2023}.

\subsubsection{CVaR-Lagrangian Objective}
\begin{equation}
\max_{\theta_{k}}\min_{\lambda_{k}\geq0}\;
  \mathbb{E}\Bigl[\sum_{t}R_{t}\Bigr]
  -\lambda_{k}\bigl(\widehat{\text{CVaR}}_{\alpha(n)}(C^{\pi_{k}})-\epsilon(n)\bigr)
\label{eq:lagrangian}
\end{equation}
Each policy has its own $\lambda_{k}$: different style preferences yield
different trajectory distributions and therefore different safety cost
distributions; a shared multiplier would distort the Pareto front.

\subsubsection{CVaR Estimation}

Collect $N$ rollouts per policy per iteration. Sort cumulative costs
$\{C^{i}\}$ descending; average the worst
$\lceil(1-\alpha(n))N\rceil$:
\begin{equation}
\widehat{\text{CVaR}}_{\alpha(n)}=
  \frac{1}{|\mathcal{T}|}\sum_{i\in\mathcal{T}}C^{i}
\end{equation}

\subsubsection{Lagrange Multiplier Update}
\begin{equation}
\lambda_{k}\leftarrow\max\!\Bigl(0,\;
  \lambda_{k}+\eta_{\lambda}\bigl(
    \widehat{\text{CVaR}}_{\alpha(n)}-\epsilon(n)\bigr)\Bigr)
\label{eq:lambda}
\end{equation}

\subsubsection{Base RL}
PPO~\cite{PPO2017} via Stable-Baselines3~\cite{SB3_2021}. Effective
per-timestep reward: $\tilde{R}_{t}=R_{t}-\lambda_{k}\cdot c_{t}$.
Policies are trained \emph{from scratch} (random initialisation) to
avoid reward bias from CaRL's scalar objective. $z_{t}$ normalisation
is bootstrapped from $N_{\text{warm}}$ IDM rollouts collected prior
to training.

\begin{algorithm}
\caption{PRISM Training}
\begin{algorithmic}[1]
\STATE Collect $N_{\text{warm}}$ IDM rollouts; compute and
       cache $\Theta$, $\epsilon$-curve, $T_{j\to i}$ values
\STATE \textbf{Stage~1}: Generate $K$ utility functions
       $\{f_{\theta_{k}}\}$
\STATE Initialise $\lambda_{k}\leftarrow0$ for all $k$
\FOR{$k=1$ to $K$}
  \STATE Initialise $\pi_{k}$ randomly
  \FOR{iteration $n=1,2,\ldots$}
    \STATE Compute $\alpha(n)$, $\epsilon(n)$ via
           Eq.~\eqref{eq:curriculum},\eqref{eq:epsilon}
    \STATE Collect $N$ rollouts; compute $R_{t}$, $c_{t}$, $C^{i}$
    \STATE Estimate $\widehat{\text{CVaR}}_{\alpha(n)}$
    \STATE PPO update with $\tilde{R}_{t}=R_{t}-\lambda_{k}c_{t}$
    \STATE Update $\lambda_{k}$ via Eq.~\eqref{eq:lambda}
    \STATE Update $z_{t}$ normalisation via EMA
  \ENDFOR
\ENDFOR
\STATE \textbf{Output}: CVaR-safe Distributional Pareto Front
       $\{\pi_{1},\ldots,\pi_{K}\}$
\end{algorithmic}
\end{algorithm}

%----------------------------------------------------------
\section{Novelty and Limitations}\label{sec:novelty}

\subsection{Contributions}

\begin{enumerate}
  \item First application of DPMORL to autonomous driving with
    domain-specific style objectives grounded in human preference
    research~\cite{Yusof2016,NatDriving2023}.
  \item CVaR safety constraint on a distributional Pareto front,
    theoretically justified by DPMORL Theorem~3~\cite{DPMORL2023}.
  \item $\alpha$-curriculum for CVaR-constrained training: progressive
    tail tightening with co-scheduled safety budget, novel to the
    safe RL literature.
  \item Principled two-tier safety cost with indicator weights derived
    from empirical causal lead times, avoiding hand-tuning.
  \item Context-aware desired speed via regime detection, addressing
    an open problem in reward design~\cite{RewardReview2024}.
\end{enumerate}

\subsection{Limitations}

Comfort and efficiency are combined into one longitudinal objective;
disentanglement is future work. Lane choice reward assumes leftmost
lane is fastest---mandatory turn scenarios are not handled in this
layer. Predicted red light violation relies on nuPlan's inferred
traffic light timing. $K\!=\!5$ discrete policies approximate the
continuous Pareto front; continuous preference
conditioning~\cite{Surmann2025} is a natural extension.

%----------------------------------------------------------
\section{Conclusion}\label{sec:conclusion}

PRISM learns a CVaR-safe Distributional Pareto Front of autonomous
driving styles, cleanly separating style optimisation from
non-negotiable safety constraints. The $\alpha$-curriculum is a
novel training mechanism applicable beyond autonomous driving to
any CVaR-constrained RL problem. We believe PRISM provides a
principled foundation for personalised, trustworthy AV systems.

%----------------------------------------------------------
\bibliographystyle{IEEEtran}
\bibliography{references}

\end{document}
